% !TEX encoding = UTF-8
%======================================================================
% 1. LỚP TÀI LIỆU VÀ ENCODING
%======================================================================
\documentclass[12pt, a4paper, oneside]{report} % onecolumn là mặc định nhưng ghi rõ cho chắc

% Hỗ trợ gõ tiếng Việt và hiển thị đúng
\usepackage[utf8]{vietnam}
%\usepackage{vntex}
\usepackage[T5]{fontenc}
%\usepackage[utf8]{inputenc}

% Thiết lập font chữ Times New Roman hoặc tương đương
\usepackage{mathptmx} % Times New Roman cho text và math
\usepackage[T5]{fontenc}
\usepackage{textcomp}


%======================================================================
% 2. CANH LỀ VÀ LAYOUT TRANG
%======================================================================
\usepackage[document]{ragged2e}

% Gói geometry để tùy chỉnh lề
\usepackage[top=3cm, bottom=3cm, left=3.5cm, right=2cm, heightrounded, headheight=15pt]{geometry}

% Gói fancyhdr để tùy chỉnh đầu trang và chân trang
\usepackage{fancyhdr}
\pagestyle{fancy}
\fancyhf{} % Xóa hết các cài đặt cũ
\lhead{\footnotesize\leftmark} % Bên trái header: tên chapter/section
\cfoot{\footnotesize\thepage} % Giữa footer: số trang
\renewcommand{\headrulewidth}{0.4pt}
\renewcommand{\footrulewidth}{0.4pt}

% Giãn dòng (1.5)
\usepackage{setspace}
\onehalfspacing

% Thụt đầu dòng cho đoạn đầu tiên của section
\usepackage{indentfirst}


%======================================================================
% 3. CÁC GÓI TOÁN HỌC
%======================================================================
\usepackage{amsmath}    % Các môi trường toán học nâng cao
\usepackage{amssymb}    % Thêm nhiều ký hiệu toán học
\usepackage{amsthm}     % Môi trường định lý, định nghĩa
\usepackage{latexsym}   % Một số ký hiệu toán học cũ
\usepackage{amsbsy}     % Ký hiệu in đậm
\numberwithin{equation}{chapter} % Đánh số công thức theo chương, ví dụ (1.1)

% Định nghĩa môi trường định lý, định nghĩa,... bằng tiếng Việt
\newtheorem{definition}{Định nghĩa}[chapter]
\newtheorem{theorem}{Định lý}[chapter]
\newtheorem{corollary}{Hệ quả}[chapter]
\newtheorem{lemma}{Bổ đề}[chapter]


%======================================================================
% 4. HÌNH ẢNH, ĐỒ HỌA VÀ MÀU SẮC
%======================================================================
\usepackage{graphicx}   % Chèn ảnh (JPG, PNG, PDF)
\usepackage{svg}        % Chèn ảnh vector SVG
\usepackage[table, svgnames]{xcolor} % Quản lý màu sắc, hỗ trợ cho bảng
\usepackage{tikz}       % Vẽ hình vector trực tiếp trong LaTeX
\graphicspath{{./images/}} % Khai báo thư mục chứa ảnh

% Cài đặt cho caption của hình ảnh, bảng
\usepackage[labelfont=bf, font=small, skip=5pt]{caption}
\usepackage{subcaption} % Hỗ trợ tạo hình/bảng con (subfigures, subtables)
\counterwithin{figure}{chapter} % Đánh số hình theo chương
\counterwithin{table}{chapter}  % Đánh số bảng theo chương


%======================================================================
% 5. BẢNG BIỂU
%======================================================================
\usepackage{booktabs}   % Tạo bảng chuyên nghiệp với các lệnh \toprule, \midrule, \bottomrule
\usepackage{tabularx}   % Bảng có thể tự động điều chỉnh độ rộng cột
\usepackage{multirow}   % Gộp ô theo hàng
\usepackage{longtable}  % Tạo bảng có thể kéo dài qua nhiều trang
\usepackage{array}      % Thêm các định dạng cột mới


%======================================================================
% 6. CHÈN MÃ NGUỒN (CODE)
%======================================================================
% Sử dụng minted cho đẹp và chuyên nghiệp, cần -shell-escape khi biên dịch
\usepackage{minted}
% Cài đặt giao diện cho minted
\setminted{
    frame=lines,      % Có khung
    framesep=2mm,     % Khoảng cách từ code đến khung
    baselinestretch=1.2,
    fontsize=\small,  % Cỡ chữ nhỏ
    linenos           % Hiện số dòng
}


%======================================================================
% 7. TÀI LIỆU THAM KHẢO VÀ HYPERLINK
%======================================================================
\usepackage{csquotes}

% biblatex là gói quản lý tài liệu tham khảo hiện đại và mạnh mẽ
\usepackage[
    backend=biber,      % Sử dụng biber (mạnh hơn bibtex)
    style=authoryear,          % Kiểu trích dẫn APA (Author-Year)
    citestyle=authoryear, % Tác giả-Năm
    sorting=nyt,        % Sắp xếp theo Name-Year-Title
    giveninits=true     % Tên tác giả chỉ viết tắt (e.g., "D. E. Knuth")
]{biblatex}
\addbibresource{bibliography.bib} % Tên file .bib của bạn

%%\DefineBibliographyStrings{vietnamese}{
%%  bibliography = {Tài liệu tham khảo},
%%  references   = {Tài liệu tham khảo},
%%}

% hyperref phải là một trong những gói được gọi cuối cùng
\usepackage[
    colorlinks=true,      % Vẫn bật colorlinks để chữ có thể có màu
    linkcolor=black,      % Màu cho link nội bộ (mục lục, section) -> ĐEN
    citecolor=black,      % Màu cho trích dẫn \cite -> ĐEN
    urlcolor=black,       % Màu cho các đường link \url -> ĐEN
    bookmarks=true,
    pdfencoding=auto,
    unicode
]{hyperref}


%======================================================================
% 8. CÁC GÓI TIỆN ÍCH KHÁC
%======================================================================
\usepackage{enumitem}   % Tùy biến môi trường liệt kê (itemize, enumerate)
\usepackage{pdfpages}   % Chèn các trang từ file PDF khác vào tài liệu
\usepackage{url}        % Hỗ trợ ngắt dòng cho các đường link dài
\usepackage[most]{tcolorbox}

% Định dạng lại tiêu đề chương, mục (section, subsection)
\usepackage{titlesec}
\titleformat{\chapter}[display]
  {\normalfont\huge\bfseries\centering}{\chaptertitlename\ \thechapter}{20pt}{\Huge}
\titlespacing*{\chapter}{0pt}{50pt}{40pt}


\usepackage[table]{xcolor}
\usepackage{float} % Đặt hình ảnh, bảng ở vị trí chính xác với [H]

\usepackage[english, vietnamese]{babel}

\usepackage{titlesec}